\section{Verification of evidence construction and D3}
\label{sec:annex-p05-pa-positive-introspection}

\begin{annexitem}{annex:p5.evidence-projections}{thm:p05-evidence-projections}{Evidence extraction and collection inside PA}
\begin{proof}
Work throughout in \(\PA\). Each function is its fixed graph output,
with existence and uniqueness from \cref{thm:p05-pa-pr-graphs}.
The list equations are those of \cref{annex:p5.checker-description};
the bounded Boolean equations are those of
\cref{thm:p05-fixed-bounded-builders}. All arithmetic inductions below
retain their other numerical variables as parameters.

Since \(Sz>0\), the head and tail equations give
\[
             Sz=C(h(Sz),t(Sz))=S(K(h(Sz),t(Sz))).
\tag{1}
\]
Successor injectivity gives \(z=K(h(Sz),t(Sz))\). The proved
identification of \(\Gamma_K\) with the original \(\mathsf P\)
therefore gives \(\mathsf P(\pi_0(z),\pi_1(z),z)\). The component
bounds in \cref{lem:p05-pa-pair-beta} give both bounds by \(z\).
Its inverse uniqueness identifies any other pair \((u,v)\) in
\(\mathsf P(u,v,z)\) with these two components. This calculation
also works at \(z=0\): both components equal zero by pairing
uniqueness and \(\mathsf P(0,0,0)\). In particular the use of
\(Sz\), rather than \(z\), is required in (1).

For the beta selector, the bounded Boolean theorem gives
\[
          b_{\mathsf B}(z,k,v)=1\leftrightarrow\mathsf B(z,k,v).
\tag{2}
\]
The internal beta lemma supplies a unique such \(v\leq z\), which
is below the search bound \(Sz\). Here is the search invariant. At
counter \(r\leq Sz\), either its flag is zero, its saved value is
zero and no counter below \(r\) satisfies (2), or its flag is one
and its saved value is the least successful counter below \(r\).
The invariant is true at zero. At the next counter, a prior flag one
is retained. With flag zero, the update either records that very counter
when its bit is one or retains the zero state when its bit is zero.
These two cases prove the invariant at \(Sr\). At \(Sz\) the first
alternative contradicts the existence just proved. The returned second
field therefore satisfies (2) and is at most \(z\); beta uniqueness
proves equality with every other supplied entry. This is an internal
induction on the specified fold, including the possibility that the saved
successful value is zero. Its flag distinguishes that success from failure.
The numerical failure default is defined even though total beta existence
proves that this particular full-bound search succeeds for all \(z,k\).

We next prove boundedness and the evidence equivalence by metainduction on
the fixed positive formation. Every newly introduced quantifier in
\(\mathcal E\) is marked bounded. Its bounds are \(z,\bar2\),
or the original bounding term, none containing its own bound variable.
The pairing matrix is atomic and beta is bounded. At each formation step,
the child evidence is already bounded. This proves boundedness of the
parent, including the bounded-universal clause; it does not exchange any
unbounded quantifier across a bounded one. The Boolean assertion for
\(e_\theta\) is now the already proved theorem for this particular
bounded formula, with precisely its free-parameter projections.

At a bounded leaf, evidence implies the original formula by right
conjunction elimination. Conversely the original formula together with
the identity \(0=0\) witnesses evidence at zero. These are ordinary
internal PA derivations, whether the leaf is positive, negated or an
implication in its marked bounded syntax.

At a conjunction, evidence supplies \(u,v\), the pairing fact and
the two child evidence facts. The two induction instances yield both
original conjuncts. Conversely those conjuncts give respective evidence
values by induction. Internal pairing totality supplies their pair code;
the component bounds give the two marked witnesses in the displayed
evidence formula. Existential introductions prove its existence. If two
conjuncts coincide, their evidence values may still differ; the pairing
clause and the proof preserve both positions.

At a disjunction, evidence supplies a tag \(a<2\), a payload \(v\)
and one of its two explicit alternatives. Equality substitution in the
first alternative gives the left child evidence, and in the second gives
the right. Apply the appropriate induction instance and disjunction
introduction, then discharge the two case assumptions. Conversely, a left
disjunct gives a payload by induction; pair zero with it. A right disjunct
gives a payload; pair one with it. The component bounds and \(0<2\),
\(1<2\) supply all evidence bounds. In either branch the actual
tagged alternative is introduced before the two bounded existential
witnesses. Disjunction elimination gives the implication from the original
disjunction to existence of evidence.

At an unbounded existential, evidence supplies the original witness and
its child evidence. The child induction instance gives the matrix at that
witness; original existential introduction gives the target. Conversely,
choose a witness for the target, then choose its child evidence using
induction. Pair these two values. Pairing bounds provide both bounds by
the resulting evidence code, so they witness exactly the displayed
evidence formula. In a marked existential, the forward direction retains
the original bound while introducing the original guarded existential.
In the reverse direction the original witness already satisfies that
bound; pairing supplies the child-evidence bound, and the same witness
and payload give the marked evidence clause. The strict and weak cases
use their original comparisons.

At a marked universal, the forward direction assumes an arbitrary value
in the original range, obtains its entry and child evidence, and applies
the child induction instance. It then discharges the range assumption
and universally generalizes that value. The other parameters and evidence
code are retained, so this proves the original universal internally.
Conversely, the original universal and the child induction instance give
\[
       \forall x_i\mathrel\diamond t(\bar x)\,
                       \exists v\,\mathcal E_\psi(\bar x,x_i,v).
\tag{3}
\]
Apply \cref{lem:p05-pa-finite-collection} to this fixed child-evidence
formula, with the other parameters retained. For a strict bound use
\(n=t(\bar x)\); for a weak bound use its proved successor-bound
version. The conclusion is precisely one code whose entry at each range
index supplies its child evidence with \(v\leq z\). This proves
the original universal evidence clause. Its strict zero case is the
finite-collection base \(z=0\); the weak zero case has the single
entry at zero. Full PA induction in that lemma licenses (3) even if the
original matrix contains unbounded existential quantifiers.

All witness and parameter substitutions in these logical derivations use
fresh auxiliary variables. In particular, existential-elimination
eigenvariables can be selected above every index in the finite major,
conclusion and remaining assumptions, so they meet the stronger
all-variable rule used by this book. When an assumption is otherwise
unused, conjoin it with the derived conclusion and project the latter;
this uses and then permits discharge of that assumption without changing
the conclusion. This gives actual nonvacuous uses for the displayed
implication and case derivations, including constant matrices.

The metainduction produces a PA proof at each fixed formation, with
universal closure in its numerical parameters. It retains the outer
parameters in every inner bound and applies finite collection only after
those parameters have been fixed. Thus a bound depending on an earlier
witness is evaluated at that witness. A shadowing binder overrides only
its own matrix parameter. Empty free-parameter lists remove only the
corresponding quantifier prefix. These observations complete all claims
of the lemma.
\end{proof}
\end{annexitem}

\begin{annexitem}{annex:p5.positive-evidence-builders}{thm:p05-positive-evidence-builders}{Every evidence clause yields the original annotated proof}
\begin{proof}
We prove finite describability, the literal root identity and the stated
safe-tree assertion together by metainduction on the fixed positive
formation. Each induction assertion is an internal theorem of PA for
all numerical tuple and evidence inputs. The outer evidence guard at
every node has already been proved equivalent to its actual bounded
evidence formula. Its failure value is zero, both before and after the
final guarded packing.

The descriptions are acyclic. The bounded leaf and evidence tests are
fixed descriptions from the preceding bounded-formula theorem. Pair
projections are compositions of already specified functions. The beta
selector is a bounded fold calling only the fixed bounded-beta test.
At a parent, every child proof description is fixed before the parent
is formed. Conjunction, disjunction and existential clauses are finite
compositions, with the displayed conditional in the disjunction clause.
The universal list uses primitive recursion with its counter last and
its old list as the sole recursive state. The resulting list is supplied
to the already specified \(\mathsf U_\diamond\). Substituting the
numerical term bound happens after the fold description has been fixed.
Every branch, including a failed child call on a failed outer input,
has a total defined numerical value. These facts permit application of
the fixed-description graph theorem without a universal evaluator.

We record the exact syntax invariant used in every clause. After its
free parameters other than a current binder \(i\) receive their
numerals, let the matrix code be \(c\). Then for each numerical
witness \(a\),
\[
              \psi^\nu(\bar m,a)=\mathsf{sub}(c,i,\nu(a)).
\tag{4}
\]
Here each side uses the child's actual free-index list; unused indices
are simply absent. The no-conflict substitution theorem proves (4) in
PA. Numeral substitutions at distinct indices commute: at a variable
they make the indicated replacement; at a nonbinding constructor use
the two child identities. At a binder for either substituted index,
that substitution stops and the other commutes into the matrix. At
every other binder both replacements are variable-free, so neither
renames the binder. Induction on the table therefore proves literal
commutation. An inner binder for \(i\) stops the operation in (4),
which proves the shadowing case. At a connective, the same equations
give the original constructor applied to its two substituted children.
The proof applies to arbitrary formula matrices, including unbounded
quantifiers; boundedness is not a premise of the substitution theorem.

At a bounded leaf, evidence gives \(\theta(\bar m)\). The fixed
bounded-formula theorem gives the safe raw tree
\(R_\theta^+(\bar m)\) with precisely root
\(\theta^\nu(\bar m)\). The evidence value is zero on successful
inputs, but the output depends only on the original parameter tuple.
For bottom the success antecedent is contradictory and no successful
zero record is asserted. All other signed literals and bounded
connectives use their already verified respective branches, including
the original guard and comparison binders.

At a conjunction, let \(a=\pi_0(z),v=\pi_1(z)\). Evidence gives
some two components satisfying \(\mathsf P\); the internal inverse
uniqueness proved above identifies those components with exactly
\(a,v\). Hence the two child evidence facts hold at these numerical
arguments. By the induction hypotheses both child raw trees are safe,
with roots \(c,d\). The specified tag-1 join has root
\(c\mathbin{\&}d\), which is exactly the parent numeral code.
Its open list is the shifted left list followed by the shifted right
list. Every leaf remains an unmarked member of \(\mathcal A_*\).
The row-join theorem gives every old and new row bit and every parent
count. Separate copies of a repeated child are inserted; no row acquires
two parents.

At a disjunction, inverse uniqueness likewise identifies the supplied
tag and payload with \(a=\pi_0(z),v=\pi_1(z)\). The first evidence
alternative gives \(a=0\) and the left child evidence. The second
gives \(a=1\) and the right child evidence. Since \(0\ne1\) in
PA, the description's first test selects exactly the first case when
that alternative holds, and its second branch uses the right evidence.
The induction hypothesis supplies the selected child safe tree. Tag 4
or 5 then gives the literal parent disjunction in its original child
order. The unused child need not have any successful proof. The outer
evidence guard rejects tags other than zero and one: being a total pair
code alone is not the success condition. Each injection retains exactly
its supplied child's axiom leaves, and the join theorem proves the tree.

At an unbounded existential, evidence and inverse uniqueness give
\[
            \mathcal E_\psi(\bar m,\pi_0(z),\pi_1(z)).
\tag{5}
\]
Let \(a=\pi_0(z)\) and \(v=\pi_1(z)\). By induction the child
tree is safe with root \(\psi^\nu(\bar m,a)\), which is the
selected substitution in (4). The new row has tag 14, parameter list
\([i,c,\nu(a)]\), one premise and conclusion \(\mathsf{Ex}(i,c)\).
The numeral is a valid variable-free term, \(c\) is a valid formula,
and its actual matrix-substitution test is (4). Thus every local
test succeeds. This rule discharges nothing and introduces no
eigenvariable; its open list is the child's shifted list. The
conclusion is the original existential binder around the original
partially substituted matrix. This includes witness zero and an
unused bound variable. In the latter case (4) still holds and the
explicit existential-introduction row still carries its witness term.

For a marked existential, the same decoding gives (5) and
\(a\mathrel\diamond t(\bar m)\). The accepted term theorem
identifies the object-language bound with the fixed numerical value
\(b=t^\#(\bar m)\). Thus \(a<N_\diamond(b)\), and the fixed
literal constructor gives a safe tree with root
\(\mathsf{sub}(G,i,\nu(a))\). Its comparison binder is the
original \(h\); a binder becoming available after closed substitution
is not newly selected. The conjunction with the child proof has root
\[
 \mathsf{sub}(G,i,\nu(a))\mathbin{\&}\mathsf{sub}(c,i,\nu(a))
                  =\mathsf{sub}(G\mathbin{\&}c,i,\nu(a)).
\tag{6}
\]
The equality is an internal no-conflict constructor equation. The
tag-14 row has parameters \([i,G\mathbin{\&}c,\nu(a)]\) and
root \(\mathsf{Ex}(i,G\mathbin{\&}c)\), exactly the original
marked existential expansion. Its two old lists contain only allowed
unmarked leaves, so their concatenation is safe. At an empty strict
range the numerical comparison and the evidence success condition
are impossible. At a weak zero bound \(a=0\) is permitted and
the actual zero-witness literal proof is used. No predecessor of
zero or fictitious witness for an empty range occurs.

Consider a marked universal. Retain all outer parameters and the
evidence \(z\). Put \(b=t^\#(\bar m)\),
\(N=N_\diamond(b)\). The fixed-term theorem and the internal order
identities identify the original range with \(k<N\). Its evidence
therefore gives, for every \(k<N\), a value \(v\leq z\) with
\(\mathsf B(z,k,v)\) and child evidence at \(k,v\). Beta
uniqueness and the selector theorem identify it with
\(\beta_*(z,k)\). We have proved inside PA
\[
 \forall k<N\,
       \mathcal E_\psi(\bar m,k,\beta_*(z,k)).
\tag{7}
\]
The child induction instance gives, for each such \(k\), a safe
tree with root \(\mathsf{sub}(c,i,\nu(k))\), by (4).

Here is the complete list-fold invariant needed to use those trees.
At counter \(r\), the list has length \(r\), and for every
\(k<r\) its entry is
\(V_\psi(\bar m,k,\beta_*(z,k))\). Under the additional
antecedent \(r\leq N\), every such entry is safe with root
\(\mathsf{sub}(c,i,\nu(k))\). At zero the list is zero and all
entry assertions are empty. At a successor append the explicitly
displayed child tree. The internal concatenation equations preserve
the old entries and place that tree at the old length. If \(Sr\leq N\),
(7) at \(r\) and the child theorem prove its safe root. The
preceding entries use the retained induction hypothesis since
\(r\leq N\). The cons-length equation increases the length by
one. This is a PA induction for the particular list description;
there is no assertion that arbitrary lists of numeral instances
already constitute a quantified derivation.

At \(r=N\) the list satisfies every hypothesis of
\cref{thm:p05-bounded-case-elimination}. The partially substituted
matrix \(c\) is a formula with free-index set contained in
\(\{i\}\); its bound term \(u=t^\nu(\bar m)\) is
variable-free; \(i\ne h\); and \(R_t(\bar m)\) is safe with
root \(\mathsf{Eq}(u,\nu(b))\). That theorem applies to this
\(c\) regardless of its unbounded existential subformulas. Its
actual construction produces the following raw roots in order:
\[
\begin{gathered}
 G_\diamond(i,h,u)\quad\hbox{as a marker-1 assumption},\\
 G_\diamond(i,h,\nu(b)),\qquad
 \mathcal D(v_i,N),\qquad c,\\
 G_\diamond(i,h,u)\Rightarrow c,\qquad
 \mathsf{All}(i,G_\diamond(i,h,u)\Rightarrow c).
\end{gathered}
\tag{8}
\]
The second root uses equality elimination with the bound-equality
tree; the third uses the actual range-exhaustion implication; the
fourth uses the case implication built from the list. The fifth
discharges the original guard, and the last universally introduces
its original binder. This is precisely the \(\mathsf U_\diamond\)
expression used in the body, not a new quantifier rule.

For completeness, the leaf conditions in this use are those proved
for that constructor. Each case minor transports equality to the
arbitrary matrix using parameters \([\nu(k),v_i,c,i]\), and uses
its equality assumption even when \(i\) is absent freely from
\(c\). Each branch filter therefore removes a surviving marked
leaf. The original guard survives through the two raw MP rows in
(8), so the implication discharge is nonvacuous. After it, every
open leaf is an unmarked member of \(\mathcal A_*\) and hence
a sentence; universal introduction passes its free-variable product.
The range-exhaustion existential eliminations choose their already
specified fresh indices above \(\max(2,i,h)\), and their
all-variable tests involve their own fixed exhaustion formulas,
major guards and remaining \(\mathcal A_*\) leaves. The new
arbitrary matrix does not occur in those earlier exhaustion rows.
The case tree adds no existential elimination. Thus an unbounded
matrix introduces no unchecked eigenvariable condition into (8).

At \(N=0\), the list is empty and the case implication is the
constructed bottom-elimination base. The strict-zero exhaustion
actually derives bottom from the guard; the last two rows of (8)
still apply. At a weak zero bound \(N=1\), the fold builds exactly
the entry at zero, using its beta-selected evidence. In dependent
inner ranges every outer numerical witness is retained as a parameter
of (7) and the list induction. Its value is supplied before computing
the inner bounding term. The literal substitution identity (4)
handles shadowing, so no witness is moved outside its own scope.

This completes the raw safe-tree metainduction. To check the sentence
assertion, use the syntax-table equations on the original formation.
At each node the remaining free indices are a subset of its declared
free-index list. Numerals have no variable occurrences. At a
quantifier the substituted matrix has at most its bound index freely;
the quantifier removes it. Conjunction and disjunction take the union
of the two child sets. Every final numeral instance consequently has
an empty free-index list and passes \(\mathsf{Sent}\), including
a formula whose parameter list was empty initially.

Unfold the actual selected characteristic description on the packed
tree. Its outer code is exactly \(N_{19}(N_{15}(R),0)\), with
the required arities and empty second field. Safety gives a nonempty
tree, all successful row bits, parent count one at every nonroot and
zero at its root, and the literal root equality. The just-proved
sentence assertion gives the sentence factor. Every final open leaf
belongs to \(\mathcal A_*\), whose nine specific codes pass
\(\mathsf{Ax}_{19}\); the bounded product of these factors is one,
including the empty-list product. Thus every actual factor of
\(\chi_{19}\) is one. Replacing its fixed graph-output notation
by \(\Gamma_{d_{19}}(p,x,\bar1)\) is its defining expansion.
It proves the selected \(\mathsf{Prf}^{19}_{\PA}\) assertion
for this literal root.

On the successful evidence antecedent the final guard returns that
packed value. Introduce its unique graph output \(p\), retain
\(\Gamma_{d_\theta}(\bar m,z,p)\), and conjoin the verified
selected-checker fact. Existential introduction in \(p\), implication
introduction and universal closure give the theorem's displayed
internal implication. All auxiliary graph-output variables are fresh
and eliminated after their unique values have been used.

For a closed \(\sigma\), its numeral substitution list is empty,
so its output root is literally \(\ulcorner\sigma\urcorner\).
Assume \(\sigma\). The evidence equivalence provides an evidence
value. The implication just proved provides a selected proof value;
discard its builder-graph conjunct and existentially introduce the
selected provability formula. Eliminate the two fresh existential
witnesses and discharge the original assumption. This proves
\(\sigma\to\mathsf{Prv}^{19}_{\PA}
(\overline{\ulcorner\sigma\urcorner})\) internally. No D3
instance or truth predicate was used to obtain it.
\end{proof}
\end{annexitem}

\begin{annexitem}{annex:p5.original-sigma-evidence}{thm:p05-original-sigma-evidence}{Signed evidence transport to every original formula constructor}
\begin{proof}
First verify the signed normal-form specification against the fixed
conversion in the hierarchy definition. Eliminating an implication
gives \(\neg\psi\lor\chi\); pushing a negative sign through it
gives \(\psi\land\neg\chi\). The two De Morgan rewrites give
the displayed conjunction and disjunction rows. Double negation
exchanges the child signs. Quantifier duality preserves the binder
and applies the negative conversion to its matrix. For a marked
quantifier, its declared bounded duality retains the original bound
and comparison; the variable is absent from its bounding term.
These rewrites never move a subformula into a different variable
scope. They therefore require no capture correction or alpha
renaming. Recursion on the fixed original syntax gives exactly the
two stated normal forms. At an atom or bottom the negative form
retains its one negation, as prescribed by pushing negations to
atomic level. Propositional abbreviations were expanded before
this procedure, under the original language convention.

The qualifying-domain condition propagates to every required signed
child. For a nonbounded connective, its normal form has exactly
the connective and children in the table, so membership in the
least positive class requires those children. At a quantifier,
the permitted bounded constructors and unbounded existential give
the listed child. An original unbounded universal with sign plus,
or unbounded existential with sign minus, would leave an unbounded
universal at that node in normal form. Such a node is not in the
stated \(\Sigma^0_1\) class. No other formation case is excluded.
If an original subformula is bounded, both signed normal forms
are bounded, so the priority clause applies for either sign.
Conversely, these rewrites retain every original unbounded-quantifier
occurrence, changing only its polarity and kind. They cannot turn
a nonbounded original subformula into a bounded normal form.
Thus outside the original bounded priority case, the normal-form
evidence has exactly the connective or quantifier clause used by
the signed constructor below.

There is one arithmetic fact needed for the bounded priority clause.
For every fixed bounded \(\eta\), PA proves
\[
 b_{\mathcal N^+\eta}(\bar m)=b_\eta(\bar m),\qquad
 b_{\mathcal N^-\eta}(\bar m)=Z(b_\eta(\bar m)).
\tag{S1}
\]
Here all bits are the previously specified fixed-formula descriptions,
with the fixed permutations of their free-input tuples. We give the
internal formation proof, so no normal-form proof code is being assumed.
At equality the positive forms coincide and the negative bit is its
zero test. At bottom the bits are zero and one. At negation the two
induction identities exchange, and \(Z(Z(b))=b\) for a zero-one
bit. At binary connectives use their exact bit operations. For
\(r,s\in\{0,1\}\), those operations and their complements are
\[
\begin{array}{c|c|c}
 \text{operation}&\text{bit}&\text{complement bit}\\\hline
 \land&rs&\mathsf{sg}(Z(r)+Z(s))\\
 \lor&\mathsf{sg}(r+s)&Z(r)Z(s)\\
 \to&I(r,s,1)&rZ(s).
\end{array}
\tag{S2}
\]
PA verifies each row by the four alternatives
\((r,s)=(0,0),(0,1),(1,0),(1,1)\), using the zero-test and
conditional equations. In the implication row its positive normal
form has bit \(\mathsf{sg}(Z(r)+s)\), which the same four
alternatives identify with \(I(r,s,1)\). Thus (S2) and the child
induction instances prove both identities at every binary node.

For a marked quantifier its two normal forms use the same bound
value and the same range. If \(r_k\) is its fixed child bit,
the necessary complement equations are
\[
\begin{aligned}
 Z\left(\prod_{k<N}r_k\right)
       &=\mathsf{sg}\left(\sum_{k<N}Z(r_k)\right),\\
 Z\left(\mathsf{sg}\left(\sum_{k<N}r_k\right)\right)
       &=\prod_{k<N}Z(r_k).
\end{aligned}
\tag{S3}
\]
They follow internally by induction on \(N\). At zero, the
product is one and the sum zero, giving respectively zero and
one on both sides. At a successor the product bit conjoins its
old bit with \(r_N\), and the sum-sign bit disjoins its old
bit with \(r_N\). The first two rows of (S2) give precisely
the two new equations. The child induction identities and the
fixed numerical bound therefore prove (S1) for each quantifier,
including dependent bounds with all outer parameters retained.
This completes the bounded formation proof of (S1).

Now prove the signed safe-tree assertions by metainduction on the
fixed original syntax, only for qualifying pairs. Every induction
assertion is an actual universally closed PA theorem for that
separately fixed signed description. Evidence inverse and beta
extraction are those already proved. The finite-description argument
is the same acyclic ordering as for the positive compiler: signs
are fixed at formation time, each child is already constructed,
and the only new repetitions are the displayed bounded list folds.
No function is evaluated at a varying formula or sign input.

At a bounded original formula, its normal-form evidence has priority
form \(z=0\land\mathcal N^\epsilon\theta\). The Boolean
equivalence for that bounded normal form and (S1) give
\(b_\theta=1\) for sign plus and \(b_\theta=0\) for sign
minus. The corresponding accepted bounded raw builder therefore
has root \(\theta^\nu\) or \(\mathsf{Neg}(\theta^\nu)\)
and is safe. This is the explicit return to the original bounded
leaf, with all original implication, negation and guard codes
retained. No proof of the normalized leaf is pasted in place of it.

At a negation, its normal form and evidence are those of the child
with exchanged sign, with the same free parameters. In the plus
case the child negative root already is the original
\(\mathsf{Neg}(c)\). In the minus case the child positive proof
has root \(c\). Contradict it with the new marker-1
\(\mathsf{Neg}(c)\) assumption. Tag 9 discharges that one leaf
and gives \(\mathsf{Neg}(\mathsf{Neg}(c))\), the exact required
negative root of the original negation. Both conclusions preserve
all child axiom leaves, and the latter discharge is nonvacuous.

For positive conjunction the evidence is a pair of its two positive
child evidence values. The two induction instances and tag 1 give
the original conjunction and the concatenated safe open list.
For negative conjunction its normal form is a disjunction of
negative children, so evidence gives tag zero with a left payload
or tag one with a right payload. Pair uniqueness identifies them
with \(a,v\). In the first case the new original conjunction
assumption is projected on the left, and contradicted with the
left negative child proof. In the second it is projected on the
right and contradicted with the right negative child proof.
The assumption survives once through that projection and
contradiction. Tag 9 removes it and gives exactly
\(\mathsf{Neg}(c\mathbin{\&}d)\). The invalid-tag cases fail
the outer evidence guard; no negative proof of the other child
is needed in a successful branch.

For positive disjunction the evidence's tag selects its positive
child and tag 4 or 5 injects that child into the original
disjunction. For negative disjunction its evidence is a pair
of negative child evidence values. Each minor contradiction uses
its own marker-1 child assumption once, with the corresponding
negative child proof. Both minors have root bottom. The tag-6
major is the marker-1 original disjunction. Its branch-specific
filters with parameters \([0,0]\) remove the respective child
assumptions nonvacuously, even when the two child formulas are
identical. They retain the major assumption. Tag 9 removes that
last marked leaf and returns the original disjunction's negation.
The only other leaves in all these trees are the children's
unmarked members of \(\mathcal A_*\).

For positive implication, its normal form is the displayed
disjunction of negative antecedent and positive consequent.
A tag-zero payload gives the negative antecedent proof. Contradict
it with the marker-1 original antecedent, eliminate bottom to the
original consequent, and use tag 7 to discharge the antecedent.
A tag-one payload gives the positive consequent proof. The
specified \(\mathsf{Keep}\) first conjoins the original
antecedent assumption with that proof and then projects the
consequent, preserving the assumption once. Tag 7 again discharges
it nonvacuously. Both branches conclude the literal
\(c\Rightarrow d\), rather than the normalized disjunction.
For negative implication its normal-form evidence supplies a
positive antecedent proof and a negative consequent proof.
Raw MP consumes the former before the original implication
assumption and gives \(d\). Contradiction with the negative
consequent gives bottom; tag 9 removes the one implication
assumption. Its root is exactly
\(\mathsf{Neg}(c\Rightarrow d)\). These are all four original
implication/branch constructions, with their actual rule parameters.

For positive unbounded existential, its evidence pair gives the
original witness \(a\) and plus-child evidence \(v\).
The child theorem gives the original matrix instance. Equation
(4) identifies its code with \(\mathsf{sub}(c,i,\nu(a))\).
The displayed tag-14 row with parameters \([i,c,\nu(a)]\)
introduces the original existential; it retains the child list.
For negative unbounded universal, its normal form is an
unbounded existential with negative matrix. Its evidence gives
the witness \(a\) and a proof of the original matrix's
negation at that witness. The marker-1 original universal is
instantiated by tag 13 with parameters \([i,c,\nu(a)]\).
The resulting positive matrix instance contradicts the negative
child proof. Tag 9 removes the universal assumption and concludes
its original negation. There is one marked leaf, and no new
eigenvariable is needed. Unused quantified variables and witness
zero still satisfy the same literal substitution identity and
the same nonvacuous assumption use.

In the two marked witness cases, the signed normal form is a
marked existential with the original bound. Its evidence pair
therefore supplies a witness satisfying that original comparison
and the required signed child evidence. The fixed literal
builder gives \(L(\bar m,a)\) with exact original guard
instance. In the positive existential, conjunction with the
plus-child tree and tag 14 give the original guarded existential
by the exact identity (6). In the negative universal, tag 13
instantiates the original marker-1 guarded universal at
\(\nu(a)\); raw MP consumes the guard proof and that
implication to give the original positive matrix instance.
Contradiction with the minus-child proof gives bottom. Tag 9
removes the universal assumption and gives its literal negation.
The successful witness condition excludes the empty strict range;
the weak zero range permits and uses witness zero.

In the two marked collection cases, the signed normal form is a
marked universal. Its evidence and beta uniqueness give child
evidence at \(k,\beta_*(z,k)\) for every \(k<N\), with the
fixed sign specified in that clause. The list induction has
length \(r\) and its prescribed child tree at each \(k<r\).
At zero all conditions are empty. At a successor the append
equations preserve each old entry and add that child's tree;
the signed child induction instance supplies its safe root
whenever \(Sr\leq N\). Thus at \(N\) the positive list has
roots \(\mathsf{sub}(c,i,\nu(k))\), and the negative list
has roots
\[
 \mathsf{Neg}(\mathsf{sub}(c,i,\nu(k)))
                    =\mathsf{sub}(\mathsf{Neg}(c),i,\nu(k)).
\tag{S4}
\]
The last equality is the nonbinding substitution equation.
The arbitrary-matrix quantified constructor applies to these
two lists. It gives respectively the original guarded universal,
or the original guarded universal with matrix \(\mathsf{Neg}(c)\).
All its row, parent-count, original-bound and original-guard
calculations were checked in (8); changing the arbitrary matrix
to \(\mathsf{Neg}(c)\) meets the same hypotheses.

For the negative marked existential, the latter quantified tree
is the displayed \(U\). Put \(E=\mathsf{Ex}(i,G\mathbin{\&}c)\).
The prescribed \(y=1+\max(2,m(E))\) is absent from every
variable occurrence of \(E\), and at least three. The minor
\(K\) therefore has literal root
\[
 G_y\mathbin{\&}c_y,\qquad
 G_y=\mathsf{sub}(G,i,v_y),\quad c_y=\mathsf{sub}(c,i,v_y).
\tag{S5}
\]
No original binder is renamed in these substitutions. Two
separate copies of that minor give \(G_y,c_y\) by their two
projections. Instantiating \(U\) gives
\(G_y\Rightarrow\mathsf{Neg}(c_y)\); raw MP with the guard
gives the negative matrix, and contradiction gives bottom.
The tag-16 minor filter, with the exact parameters
\([i,G\mathbin{\&}c,y,1]\), removes both surviving marker-2
copies of (S5). It is nonvacuous. The major is the original
marker-1 \(E\), and the result's open list is the filtered
minor list followed by the major list. Apart from that major,
all remaining leaves belong to \(\mathcal A_*\) and have
indices only zero and one. The conclusion bottom has no variables;
the major has none at \(y\) by its choice. Hence every actual
factor of the all-variable freshness test is one. Tag 9 removes
the sole marked major and concludes exactly
\(\mathsf{Neg}(E)\). This verifies the fourth marked case,
including unbounded quantifiers anywhere inside its matrix.

At a strict zero bound both collection lists have length zero.
The quantified constructor still uses its bottom case and the
strict-zero exhaustion proof. Thus positive universal and
negative existential succeed from their empty universal evidence
without obtaining an index witness. The negative existential
still has its two-copy minor and both nonvacuous discharges just
checked. At a weak zero bound there is exactly one collected
case. Bound-variable shadowing and dependent inner bounds satisfy
(4) and retain all outer numerical parameters in the list
induction, as in the positive construction. These statements
cover every zero, empty and binder case in the signed clauses.

For each row just described, the actual template specifies its
rule tag, exact root, ordered premises, parameter list and marker.
The template theorem gives successful copied row bits, one
parent at every old root and nonroot, zero at the new root,
and the exact shifted open-list update. The explicit checks
above supply every discharge and quantifier side condition;
all remaining leaves are unmarked members of \(\mathcal A_*\).
This proves the signed safe-tree assertions internally at every
formation step. Its root-code calculation uses only the original
connective and binder constructors and the selected substitution
equations. The normal-form code enters only the evidence test
and the choice of child evidence shape.

The original formula's numeral instance is a sentence by its
free-index substitution invariant; its negation is a sentence
by the same syntax-table equation. The safe invariant therefore
gives all actual tag-19 checker factors on packing, exactly as
in the positive proof. The outer evidence guard returns the
packed record on successful inputs and zero on all failed inputs.
Introducing the unique output graph of
\(\widehat d_\theta^\epsilon\) gives its selected-checker
conjunction with the respective original signed root.

Finally fix any closed \(\Sigma^0_1\) sentence \(\sigma\)
in its original presentation. Its definition of that class says
precisely that \(\mathcal N^+\sigma\) is a qualifying positive
formation. Its existing \(\mathsf{Ev}_\sigma\) is the evidence
of that very fixed normal form, including the inherited auxiliary
choices. Thus the plus-sign theorem applies to that exact evidence.
The parameter tuple is empty; the output root is
\(\ulcorner\sigma\urcorner\), by the original-syntax invariant,
so it yields the displayed unary-builder implication in the body.
Internal evidence existence for \(\sigma\), followed by that
implication and existential introduction in its proof value,
gives \(\sigma\to\mathsf{Prv}^{19}_{\PA}
(\overline{\ulcorner\sigma\urcorner})\). The ordinary PA
evidence-existence theorem provides the logical normal-form
equivalence; the signed record construction above supplies the
additional return to the original quoted code. Neither an
arbitrarily chosen equivalent formula nor D3 is used in that
record construction.
\end{proof}
\end{annexitem}

\begin{annexitem}{annex:p5.literal-graph-evidence}{thm:p05-literal-graph-evidence}{Literal graph formation, recursion traces and the outer sentence}
\begin{proof}
Inspect the actual graph definition of
\cref{def:p05-pa-pr-graph}. Its initial clauses are equalities
\(y=0,y=Sx,y=x_i\). The original pairing formula is the single
equality
\[
              (a+b)\times S(a+b)+(b+b)=w+w.
\tag{9}
\]
The original beta formula is
\[
 \exists t<Sw\,\exists m<Sw\,\exists q<St\,
  \bigl(\mathsf P(t,m,w)\land t=d\times q+v
                   \land\exists u<d\,(v+Su=d)\bigr),
 \qquad d=S((Si)\times m).
\tag{10}
\]
All its quantifiers are marked bounded. After expanding (9) and its
term abbreviation, (10) contains only equalities, conjunctions and
marked bounded existentials. It is already in the declared marked
negation normal form. The auxiliary comparison witnesses in a marked
guard are retained when that abbreviation is later expanded to its
actual tree code. The hierarchy conversion acts on the marked
presentation, as required by its definition.

Metainduct now on a fixed graph description. An initial graph has the
property by its equality clause. A composition adds only its specified
unbounded existential variables and conjunctions of child graphs.
The property is preserved. A recursion adds its outer trace
existential, its initial-entry existential, a marked bounded universal
whose matrix has two existentials and conjunctions, and its last beta
conjunct. Its child graphs and all beta occurrences have the property
by induction and (10). These are exactly the positive-formation
constructors. They introduce no negation and no unmarked implication.
The implication in the fully expanded code of a marked universal
continues to carry its declared bounded mark for the hierarchy
conversion. It is not first stripped of that mark and then normalized.
Thus the fixed marked negation-normal-form conversion changes none of
these graph nodes.

Their graph-bound variables were selected fresh before substitution,
so the conversion needs no freshening at a subsequent numeral input.
A variable-free numeral cannot cause capture, and numeral substitution
commutes with every retained constructor by the no-conflict theorem.
Hence every fixed numeral instance has the same property, with its
literal original bound indices. The evidence construction therefore
applies precisely its original clauses and priority to this same
formation. With the same auxiliary fresh choices it is literally
\(\mathsf{Ev}\). No equivalence of two separately coded graph
formulas is needed in this argument.

We spell out the composition and trace uses of the builder to account
for every evidence position. In a composition, its original formula
has its fixed consecutive binders \(z_1,\ldots,z_m\) and its
originally associated conjunction of the \(h_j\) graphs and the
final \(g\) graph. At each outside existential evidence node,
\(\pi_0\) gives the numerical value of that particular intermediate
variable, and \(\pi_1\) gives the remaining evidence. Repeating
these specified projections \(m\) times supplies all intermediate
values and the conjunction evidence. Each original binary
conjunction node is decoded by its two projections, and the
corresponding two child proofs are joined in that same association.
The final \(g\) graph receives exactly those intermediate numerical
values in its declared argument positions. Then the \(m\) actual
existential-introduction rows wrap the conjunction from the innermost
binder to the outermost, with these values' numeral terms. Each row
has its own partially substituted original matrix as its parameter.
Equation (4) verifies each literal root. When \(m=0\), the original
definition inserts no intermediate quantifier or conjunct; the
builder uses the remaining \(\Gamma_g\) formation directly. It
does not introduce a bottom or a spurious empty conjunction.

For a recursion, keep its original fresh indices and abbreviate the
three original conjuncts, only in this explanation, by
\[
\begin{aligned}
 A(w)&=\exists u(\mathsf B(w,0,u)\land\Gamma_g(\bar x,u)),\\
 H(w,k)&=\exists u\exists v
   (\mathsf B(w,k,u)\land\mathsf B(w,Sk,v)
                            \land\Gamma_h(\bar x,k,u,v)),\\
 Z(w)&=\mathsf B(w,n,y).
\end{aligned}
\tag{11}
\]
The literal formula is its original \(\exists w\) around the
original association of \(A(w)\), \(\forall k<n\,H(w,k)\),
and \(Z(w)\). If a whole subformula qualifies for the bounded
priority clause, the bounded builder supplies that exact subformula
at once. Otherwise the following clauses give its construction.

The outside evidence pair first supplies a numerical trace code
\(w_0\) and evidence for those three conjuncts at that trace code.
The conjunction projections supply separate evidence for all three
original positions. For \(A(w_0)\), its existential evidence gives
an initial value \(u_0\) and a conjunction payload. The beta leaf
builder proves precisely \(\mathsf B(\overline{w_0},0,
\overline{u_0})\); the fixed base-graph builder proves its original
\(\Gamma_g\) instance. Join them and introduce the original
initial-entry binder using \(\nu(u_0)\).

For the middle conjunct let its evidence code be \(s\) and its
numerical bound be the input \(n_0\). Its proof-list entry at
each \(k<n_0\) uses the numerical evidence
\(\beta_*(s,k)\). This selector reads the evidence sequence,
which is distinct from the original computation trace \(w_0\).
At the \(H\)-evidence node, two consecutive pair projections
supply the original step values \(u_k,v_k\) and their remaining
conjunction payload. Its three child proofs have precisely the
roots
\[
\begin{gathered}
 \mathsf B(\overline{w_0},\bar k,\overline{u_k}),\qquad
 \mathsf B(\overline{w_0},S\bar k,\overline{v_k}),\\
 \Gamma_h(\overline{\bar x},\bar k,
                              \overline{u_k},\overline{v_k}).
\end{gathered}
\tag{12}
\]
In particular the second argument in the second beta formula is the
original term \(S\bar k\), not a silently replaced numeral term.
The fixed-term transport within the bounded builder and the original
matrix substitutions retain precisely that term code. Join the
three proofs in their original association. Introduce the original
\(v\) binder and then the original \(u\) binder, using their
numerals and their exact partially substituted matrices. The result
is the list entry with root the original \(H(w_0,\bar k)\).
The universal builder then converts this complete list through (8)
to the original \(\forall k<\bar n_0\,H(w_0,k)\), retaining
its original comparison binder. At \(n_0=0\), this step uses
the zero-case exhaustion construction, without reading any step
evidence or asserting any \(u_k,v_k\).

The last conjunct's bounded builder proves the exact original
\(\mathsf B(\overline{w_0},\bar n_0,\bar y_0)\). The
original input/output arguments may be fixed terms rather than
variables; their retained term codes and fixed-term transports give
the same statement with those terms. Join the initial, universal
and last proofs in their original association and introduce the
original trace binder with \(\nu(w_0)\). This completes every
conjunct of the literal recursion formula. Neither validity of the
numerical trace alone nor a finite list of step facts was used as a
substitute for these rows.

Apply this finite formation argument to the actual expanded
\(d_{19}\). Its many library calls are compositions and recursions
already fixed in that description. The argument therefore covers
all their graph clauses, including branches computing a zero value.
Numerical negated tests and conditionals in the checker specification
have been expanded into their zero-test, composition and recursion
descriptions before \(\Gamma\) is formed. They consequently do not
insert negations or implications of child graph formulas into the
selected \(\Gamma_{d_{19}}\).
For evidence of the selected proof relation its output argument is
the fixed term \(\bar1\), and its input formula-code argument is
\(\overline{n_\varphi}\). The current construction does not
replace the relation by a separately chosen characteristic formula.

Finally inspect the original provability binder. It differs from
the free formula-code index \(j\). Since \(\nu(n_\varphi)\)
has no variables, the selected substitution equation gives
\[
 \mathsf{sub}(\mathsf{Ex}(i,g_{19}),j,\nu(n_\varphi))
       =\mathsf{Ex}(i,\mathsf{sub}(g_{19},j,\nu(n_\varphi)))
       =c_\varphi.
\tag{13}
\]
For this fixed closed sentence, its evidence is literally
\[
 \exists a\leq z\,\exists v\leq z\,
    (\mathsf P(a,v,z)\land\mathcal E_{\gamma_\varphi}(a,v)),
\tag{14}
\]
with the original evidence auxiliary choices. Pair inverse uniqueness
turns (14) into the child evidence at
\(a=\pi_0(z),v=\pi_1(z)\). The child theorem supplies a safe
raw tree whose root is exactly
\(\mathsf{sub}(g_\varphi,i,\nu(a))\). The final row of
\(d^3_\varphi\) is tag 14, parameter list
\([i,g_\varphi,\nu(a)]\), root
\(\mathsf{Ex}(i,g_\varphi)=c_\varphi\). It discharges nothing
and keeps only the child's allowed unmarked leaves. Packing and its
success guard give every selected checker factor by the preceding
theorem. Introducing the unique builder output gives precisely the
displayed \(\Gamma_{d^3_\varphi}\) conjunction and its universal
implication in \(z\).

The root arithmetic argument \(\overline{c_\varphi}\) is the
numeral of this literal formula code. Formula-code operations at
the fixed standard value \(n_\varphi\) have their specific
finite graph computations, so PA proves their values in (13).
This justifies replacing the fixed code-expression abbreviations
by that numeral in the final assertion. It is a calculation of
syntax constants, not an inference from extensional agreement of
proof predicates. No normal-form or quote-code transport is omitted:
the graph's marked normal-form conversion is the identity, and
(13) is the exact selected substitution identity.
\end{proof}
\end{annexitem}

\begin{annexitem}{annex:p5.direct-pa-d3}{thm:p05-direct-pa-d3}{The internal D3 instance for the same literal formula}
\begin{proof}
Fix an arithmetic sentence \(\varphi\), and retain precisely
\(n_\varphi,c_\varphi,\sigma_\varphi,d^3_\varphi\) from
\cref{def:p05-literal-introspection-instance}. The sentence
\(\varphi\) itself need not be \(\Sigma^0_1\). Its selected
provability sentence \(\sigma_\varphi\) has the original
positive graph formation just verified. Therefore PA proves both
\[
\begin{aligned}
 \sigma_\varphi&\to\exists z\,
                       \mathsf{Ev}_{\sigma_\varphi}(z),\\
 \forall z\,\bigl(\mathsf{Ev}_{\sigma_\varphi}(z)&\to
   \exists p\,(\Gamma_{d^3_\varphi}(z,p)\land
                 \mathsf{Prf}^{19}_{\PA}(p,\overline{c_\varphi}))\bigr).
\end{aligned}
\tag{15}
\]
The first is the internal evidence-existence theorem for this same
formula. The second is its constructed original-root proof invariant.

Here are the remaining logical steps in PA. Assume
\(\sigma_\varphi\). Apply the first line of (15). Eliminate
its existential with a variable fresh from both fixed sentences
and all open assumptions. Under that evidence assumption, instantiate
the second line at the fresh variable and use implication
elimination. Eliminate the resulting existential proof value with
another such fresh variable. Conjunction elimination gives
\(\mathsf{Prf}^{19}_{\PA}(p,\overline{c_\varphi})\).
Introduce its proof existential, choosing the original bound name
of \(\mathsf{Prv}^{19}_{\PA}\) and the current fresh witness
term. The result is literally
\(\mathsf{Prv}^{19}_{\PA}(\overline{c_\varphi})\).
Its original graph expansion and the premise agree by
capture-avoiding substitution; auxiliary graph binders had already
been fixed fresh. If the working existential witness has a different
name, the explicit alpha rule identifies only its temporary
renaming with the original selected matrix, with equal binding
positions and unchanged free indices.

The last sentence contains neither fresh witness variable. Choose
both witnesses above every index occurring in the two fixed PA
derivations in (15), their open axiom leaves, the premise and
conclusion; these are finite sets. Thus the all-variable tests for
both existential eliminations hold. Each minor assumption was used
to obtain the displayed instance or conjunction, so each discharge
is nonvacuous. Eliminate the two temporary witnesses. The original
\(\sigma_\varphi\) assumption remains used in the first
implication elimination; implication introduction discharges it
and gives
\[
       \PA\derives\sigma_\varphi\to
                 \mathsf{Prv}^{19}_{\PA}(\overline{c_\varphi}).
\tag{16}
\]

By the literal code identity (13),
\(c_\varphi=\ulcorner
\mathsf{Prv}^{19}_{\PA}(\overline{\ulcorner\varphi\urcorner})
\urcorner\). Expanding the two names in (16) therefore gives
exactly the D3 formula in the body. This last replacement changes
no predicate, numeral convention or graph binder. The output proof
constructed from evidence proves the original provability sentence;
it is not a proof of its own checker assertion obtained by assuming
D3. Both input and output provability predicates are the same
\(\exists p\,\Gamma_{d_{19}}(p,x,\bar1)\) used in the accepted
D1 and D2 construction.

All existence and induction assertions used above are PA theorems
for separately fixed descriptions and formulas. Neither consistency
nor soundness of PA was a premise. The fixed sentence parameter can
be replaced externally by any other arithmetic sentence to obtain
its corresponding instance and its separately specified builder.
This proves the stated scheme without quantifying inside PA over
all primitive recursive descriptions.
\end{proof}
\end{annexitem}
